\chapter{CHAPTER 5: SYSTEM
IMPLEMENTATION}\label{chapter-5-system-implementation}

\section{5.1 Chapter Overview}\label{chapter-overview}

This chapter details the technical implementation of the SentinelAQ
system. To authentically demonstrate the engineering lifecycle and
analytical problem-solving required for this study, the implementation
is presented as an iterative evolution through three distinct versions.
Section 5.2 documents Version 1 (v1), which attempted cutting-edge
satellite fusion but encountered catastrophic environmental constraints.
Section 5.3 describes Version 2 (v2), a pivot to resolve data continuity
that highlighted the limitations of spatially blind modeling. Section
5.4 details Version 3 (v3), the final, robust architecture integrating
mathematical feature engineering, Explainable AI (SHAP), and reliable
meteorological proxies. Finally, Section 5.5 outlines the integration of
the backend pipeline with the React Native mobile application.

\section{5.2 Version 1 (v1): Satellite Fusion and the Cloud Occlusion
Problem}\label{version-1-v1-satellite-fusion-and-the-cloud-occlusion-problem}

\subsection{5.2.1 Initial Concept}\label{initial-concept}

The initial technical hypothesis (v1) aimed to construct a highly
sophisticated multi-modal architecture. The objective was to fuse
high-frequency ground telemetry (PurpleAir PM2.5) with daily chemical
vectors (Tropospheric NO2) extracted from the Sentinel-5P TROPOMI
satellite. The theoretical logic was sound: use NO2 as a spatial proxy
to detect incoming traffic pollution before it reached the ground
sensors in Colombo and Kandy.

\subsection{5.2.2 Implementation and
Failure}\label{implementation-and-failure}

A data extraction script was deployed within the Google Earth Engine
(GEE) cloud architecture, executing spatial reduction functions over the
exact coordinates of the target cities. However, during the Python
preprocessing phase, massive anomalies emerged. Sri Lanka's tropical
monsoons, combined with Kandy's deep valley topography, resulted in
persistent, dense cloud cover. Because optical satellite sensors cannot
penetrate clouds, the dataset suffered from severe data loss---exceeding
40\% of the total observational days.

When standard statistical imputation techniques (such as K-Nearest
Neighbors) were applied to fill these massive gaps, they introduced
synthetic noise that destroyed the underlying temporal signal.
Consequently, the initial XGBoost and LSTM training runs failed
catastrophically, yielding excessively high Root Mean Squared Errors
(RMSE) during short-term forecasting.

\subsection{5.2.3 Critical Analysis and
Pivot}\label{critical-analysis-and-pivot}

The v1 failure yielded a crucial insight: while multi-modal satellite
fusion is highly effective in temperate, cloud-free environments (as
seen in European literature), it is fundamentally unviable for
high-frequency, real-time public health alerting in tropical
microclimates. A structural pivot was required to ensure system
reliability.

\section{5.3 Version 2 (v2): Ground-Only Sensing and Feature
Scarcity}\label{version-2-v2-ground-only-sensing-and-feature-scarcity}

\subsection{5.3.1 The Pivot for Data
Continuity}\label{the-pivot-for-data-continuity}

To resolve the catastrophic data gaps of v1, the architecture was
pivoted in Version 2 (v2) to rely exclusively on ground-based sensing. A
Python script utilizing the \texttt{requests} library was engineered to
execute HTTP GET operations against the PurpleAir REST API, fetching
hourly PM2.5 averages for Colombo (ID: 29677) and Kandy (ID: 12451).
This pivot successfully reduced missing data from \textgreater40\% to
\textless2\%, which was easily managed via linear interpolation using
the \texttt{pandas} library.

\subsection{5.3.2 Implementation and
Limitations}\label{implementation-and-limitations}

With a continuous dataset secured, an initial baseline XGBoost model was
trained. The feature vector was intentionally simplistic, relying solely
on historical PM2.5 lags (e.g., pollution levels at t-1, t-2, and t-24
hours).

While this solved the data continuity crisis, the resulting model was
functionally ``spatially blind.'' Evaluation metrics revealed that the
model could predict the immediate t+1 hour trend reasonably well due to
simple autoregression (momentum). However, performance degraded
exponentially at the 24-hour and 48-hour horizons. Without
meteorological context (wind, temperature), the algorithm could not
mathematically anticipate when a locally trapped pollution cloud in
Kandy would disperse, or when transboundary smog from the ocean would
strike Colombo.

\subsection{5.3.4 Critical Analysis}\label{critical-analysis}

The v2 iteration proved that while high-frequency data continuity is a
prerequisite for machine learning, it is insufficient on its own. To
achieve accuracy across extended 48-hour horizons, the architecture
required contextual environmental variables that did not suffer from the
occlusion problems of satellite imagery.

\section{5.4 Version 3 (v3): The Final Robust
Architecture}\label{version-3-v3-the-final-robust-architecture}

\subsection{5.4.1 Integrating Reliable Meteorological
Proxies}\label{integrating-reliable-meteorological-proxies}

Version 3 (v3) represents the finalized SentinelAQ architecture. To
provide the necessary spatial and environmental context without the
fragility of satellite data, the Open-Meteo API was integrated. Python
extraction scripts were updated to concurrently fetch hourly Wind Speed,
Wind Direction, Precipitation, Temperature, and Humidity.

\subsection{5.4.2 Advanced Feature Engineering (The 48-Feature
Vector)}\label{advanced-feature-engineering-the-48-feature-vector}

A sophisticated data engineering pipeline was developed using
\texttt{pandas} and \texttt{numpy}. Recognizing the failures of v2, the
simple lag vector was expanded into a highly complex 48-dimensional
matrix. Crucial engineered variables included: * \textbf{Rolling
Statistical Means:} Moving averages calculated over 3h, 6h, 12h, and 24h
windows. This explicitly programmed the concept of ``pollution
accumulation'' into the tabular data, drastically improving the model's
ability to predict smog events in the Kandy valley. * \textbf{Cyclical
Time Encoding:} Machine learning algorithms inherently treat time as
linear (e.g., Hour 23 is mathematically ``far'' from Hour 0, despite
them being contiguous). To resolve this, trigonometric sine and cosine
transformations were applied to the `Hour of Day' and `Month of Year'
variables, explicitly mapping the cyclical nature of daily traffic
rushes and seasonal monsoons into a format the algorithms could natively
understand.

\subsection{5.4.3 Multi-Horizon Training
Pipeline}\label{multi-horizon-training-pipeline}

The enriched v3 dataset was partitioned via a strict chronological split
(2022--2024 for training, 2025 for testing) to prevent look-ahead bias.
Three distinct architectures were trained across five horizons (1h, 6h,
12h, 24h, 48h): 1. \textbf{XGBoost:} Hyperparameter tuned using
\texttt{GridSearchCV} to optimize tree depth (\texttt{max\_depth=6}) and
learning rate, effectively modeling the non-linear feature interactions.
2. \textbf{LSTM Encoder-Decoder:} Constructed via
\texttt{tensorflow/keras}, utilizing a 24-hour sequential lookback
window with dropout layers to prevent the overfitting observed in early
tests. 3. \textbf{SARIMAX:} Implemented via the \texttt{statsmodels}
library to serve as the rigorous statistical baseline.

\subsection{5.4.4 Explainable AI (SHAP)
Integration}\label{explainable-ai-shap-integration}

To completely deconstruct the ``black box'' nature of the v3 ML models,
the \texttt{shap} Python library was integrated. The
\texttt{shap.TreeExplainer} was applied directly to the optimal XGBoost
model. This algorithm mathematically calculates Shapley values for every
single hourly forecast, isolating the precise marginal contribution of
features (e.g., proving that high wind speed negatively impacted the
PM2.5 prediction, while a high 24-hour rolling mean positively impacted
it).

\section{5.5 Mobile Application and Cloud
Integration}\label{mobile-application-and-cloud-integration}

The final phase of implementation bridged the Python machine learning
backend with an intuitive user interface. * \textbf{Firebase
Middleware:} The Python inference script was containerized and
scheduled. Upon generating a multi-horizon forecast and its accompanying
SHAP explanations, the backend structures a JSON payload and executes a
write operation to a Google Firebase Firestore NoSQL database via the
\texttt{firebase-admin} SDK. * \textbf{React Native Frontend:} A mobile
application was developed to consume this data in real-time. Upon
launch, the app fetches the latest Firestore document, rendering the
current AQI via a dynamic color-coded dial (Green to Red based on US EPA
standards) and plotting the 48-hour forecast on an interactive chart. *
\textbf{Automated Alerting:} The system includes a Firebase Cloud
Messaging (FCM) trigger. When the backend predicts an AQI spike that
breaches a user-defined threshold within the next 48 hours, a targeted
push notification is dispatched to the mobile device, fulfilling the
system's mandate as a proactive public health tool.
